\documentclass{beamer}
\usetheme{Madrid}
\usecolortheme{beaver}
\usepackage{amsmath}

\title{Collaborative Writing with CRDTs}
\subtitle{A short talk}
\author{Your Name}
\institute{University of Chittagong}
\date{\today}

\begin{document}

\begin{frame}
  \titlepage
\end{frame}

\begin{frame}{Outline}
  \tableofcontents
\end{frame}

\section{Problem}
\begin{frame}{Why is collaborative editing hard?}
  \begin{itemize}
    \item Two people type at the same place at the same time
    \item Someone edits offline on a bus
    \item The network delivers messages out of order
  \end{itemize}
  \begin{block}{Goal}
    Everyone ends up with the same document, without losing anyone's work.
  \end{block}
\end{frame}

\section{Idea}
\begin{frame}{Give every character an identity}
  \begin{columns}
    \column{0.5\textwidth}
      Each character remembers its left and right neighbours when it was typed.
    \column{0.5\textwidth}
      \[
        c = (\text{author}, \text{counter}, \text{origin}_L, \text{origin}_R)
      \]
  \end{columns}
  \begin{alertblock}{Result}
    Concurrent insertions are ordered deterministically and never interleave.
  \end{alertblock}
\end{frame}

\begin{frame}{Thank you}
  \centering \Large Questions?
\end{frame}

\end{document}
